% !TEX root = ../bb_AiRa_Kappaleet.tex

%% HARRIS: Chapter 10

% ö = \%c3\%b6
% ä = \%C3\%A4

\xdef\sectiontitle{\IIsectiontitle} %aiheuttama
\TOCrefs

%%%%%%%%%%%%%%%
\begin{frame}[label=2_4_a]%[label=2_3_TOC1]
\frametitle{\texorpdfstring{\culine[\sectiontitleulinecolor]{\sectiontitle}}{\sectiontitle} }
\label{\TOCname}

\vspace{\chapterTOCvksip}
\begin{itemize}[leftmargin=0.5cm,itemsep=3mm]
    \item[2.1]<1-> \hyperlink{2_1_a}{\IIaaSubsectiontitle}
    \item[2.2]<1-> \hyperlink{2_2_a}{\IIaSubsectiontitle}
    \item[2.3]<1-> \hyperlink{2_3_a}{\IIbSubsectiontitle}	
    \item[2.4]<1-> \hyperlink{2_4_a}{\CDAlert[blue]{\IIcSubsectiontitle}}	
    \item[2.5]<1-> \hyperlink{2_5_a}{\IIdSubsectiontitle}
    \item[2.6]<1-> \hyperlink{2_6_a}{\IIeSubsectiontitle}
    \item[2.7]<1-> \hyperlink{2_7_a}{\IIfSubsectiontitle}
    \item[2.8]<1-> \hyperlink{2_8_a}{\IIgSubsectiontitle}
    \item[2.9]<1-> \hyperlink{2_9_a}{\IIhSubsectiontitle}
    \item[2.10]<1-> \hyperlink{2_10_a}{\IIiSubsectiontitle}
\end{itemize}


%\endofslide 
\end{frame}
%%%%%%%%%%%%%%%

\xdef\sectiontitle{\IIcSubsectiontitle} 

%%%%%%%%%%%%%%%
\begin{frame}
\frametitle{\texorpdfstring{\culine[\sectiontitleulinecolor]{\sectiontitle}}{\sectiontitle} }
\framesubtitle{Ympäröivien atomien vaikutus}
\appear{}

\begin{itemize}
	\item Aineet \appear{koostuvat \CDAlert[red]{erittäin suuresta määrästä atomeita}}

\item Aineella on \CDAlert[black]{kolme fyysistä} olomuotoa \appear{tai \CDAlert[black]{faasia}:} \appear{\CDAlert[violet]{kaasut}}\appear{, \CDAlert[blue]{nesteet}}\appear{, ja \CDAlert[red]{kiinteät aineet}} \appear{(4:s olomuoto, \CDAlert[fuchsia]{plasma}, jätetään tarkastelematta)}

\item Kaasuissa molekyylien keskimääräiset etäisyydet  \appear{ovat paljon suurempia kuin molekyylien koot} 
\implies{ \CDAlert[violet]{Molekyylit säilyttävät yksilöllisyytensä} }

\item \CDAlert[red]{Kiinteissä aineissa atomit} (tai molekyylit) ovat \CDAlert[red]{tiiviisti pakkautuneita}\appear{, ja ovat sähkömagneettisten voimien \CDAlert[red]{paikoilleen} vangitsemia}\appear{. Näiden voimien suuruusluokat ovat samaa luokkaa kuin molekyylisidosten väliset sidosenergiat} 

\item Kiinteissä aineissa atomit eivät säilytä yksilöllisyyttänsä\appear{, vaan niiden \CDAlert[red]{ominaisuudet riippuvat ympäröivistä atomeista}}

\end{itemize}
% two last form condensed matter .. add a note
%Three phases of water (Richard Feynman, 1964 Lectures in Physics, Voll)

\endofslide 
%\endofsection
\end{frame}
%%%%%%%%%%%%%%%


%%%%%%%%%%%%%%%
\begin{frame}
\frametitle{\texorpdfstring{\culine[\sectiontitleulinecolor]{\sectiontitle}}{\sectiontitle} }
\framesubtitle{Tiiviin vs. kiinteän aineen fysiikka}
\appear{}

\begin{itemize}
\item \CDAlert[blue]{Nesteet} ovat kaasujen ja kiinteiden aineiden välimaastossa 
\item[] \begin{itemize}
	\item Molekyylit muodostavat \CDAlert[blue]{väliaikaisia lyhyen matkan} sidoksia \appear{jotka jatkuvasti hajoavat ja uudelleenmuodostuvat} \appear{(kineettisen energian ajamina)}
\end{itemize}

\item Moderni fysiikan osa-alue \CDAlert[blue]{tiiviin aineen fysiikka} \appear{\Alert[blue]{pitää sisällään sekä \CDAlert[blue]{kiinteiden aineiden että nesteiden}} tutkimisen}

\item Kuten nimestä voi päätellä\appear{, \CDAlert[red]{kiinteän olomuodon fysiikka}} \appear{\Alert[red]{keskittyy 'vain' kiinteisiin aineisiin} ja niiden ominaisuuksiin}

\end{itemize}
%

\endofslide 
%\endofsection
\end{frame}
%%%%%%%%%%%%%%%

%%%%%%%%%%%%%%%
\begin{frame}
\frametitle{\texorpdfstring{\culine[\sectiontitleulinecolor]{\sectiontitle}}{\sectiontitle} }
\framesubtitle{Kiteytyminen}
\appear{}

\begin{itemize}[itemsep=1.8mm]
	\item \CDAlert[red]{Suurin osa kiinteistä aineista on kiteisiä}\appear{, joissa atomit, ionit tai molekyylit} \appear{muodostavat säännöllisiä}\appear{, toistuvia kolmiulotteisia rakenteita}\appear{, muodostaen kidehilan} 
	\implies{ Kiinteillä aineilla on \CDAlert[red]{pitkän matkan järjestystä/säännöllisyyttä}}
	
\item \CDAlert[tunired]{Amorfisilla kiinteillä aineilla} \appear{(esim. lasit, kivet ja kumi)} \appear{niiden rakenneosaset eivät ole säännöllisesti järjestyneet}\appear{, eli atomit/molekyylit eivät ole periodisesti järjestyneet}

\item Jos nestettä jäähdytetään hitaasti\appear{, molekyylien kineettinen energia pienenee ja aineen jäätyessä molekyylit järjestäytyvät säännölliseksi kiteiseksi aineeksi}\appear{, suurin määrä sidoksia muodostuu, mikä minimoi potentiaalienergian}

\item Jos \CDAlert[tunired]{neste jäähdytetään nopeasti}\appear{, niin että systeemin sisäenergia häviää nopeammin kuin mitä molekyylien järjestäytyminen vie aikaa}\appear{, jäätyneestä aineesta tulee \CDAlert[tunired]{amorfinen}}\appear{, 'valokuva' nesteestä}
\end{itemize}

\endofslide 
%\endofsection
\end{frame}
%%%%%%%%%%%%%%%

%%%%%%%%%%%%%%%
\begin{frame}
\frametitle{\texorpdfstring{\culine[\sectiontitleulinecolor]{\sectiontitle}}{\sectiontitle} }
\framesubtitle{Kiteytyminen}
\appear{}

\begin{itemize}[itemsep=1.9mm]
	\item Suurin osa kiinteistä aineista on \CDAlert[red]{monikiteisiä}\appear{, eli ne koostuvat useista yksittäisistä kiteistä}

\item Yksittäisten kiteiden koot voivat olla millimetrien (tai senttien) suuruusluokkaa 

\item Kiteiden koot riippuvat esim. muodostumisprosessin lämpötilahistoriasta (nopea vs. hidas jäähtyminen) 

\item Yksittäinen kide voidaan muodostaa pienimmän rakenneosasen\appear{, \CDAlert[red]{ yksikkökopin avulla}}\appear{, jota toistamalla koko kide voidaan muodostaa}\appear{, \CDAlert[tunired]{määrittäen kiteen symmetrian ja säännöllisyyden}}

\item Yksikkökopin tyyppi riippuu \appear{atomien}\appear{, ionien}\appear{ tai molekyylien} \appear{\CDAlert[red]{välisten sidosten tyypistä}}
\end{itemize}

\endofslide 
%\endofsection
\end{frame}
%%%%%%%%%%%%%%%

%%%%%%%%%%%%%%%
\begin{frame}
\frametitle{\texorpdfstring{\culine[\sectiontitleulinecolor]{\sectiontitle}}{\sectiontitle} }
\framesubtitle{Kiteiden epämuodostumat, jännitys ja venymä}
\appear{}

%Crystal defects
\begin{itemize}
	\item Täydellisessä kiteessä \appear{jokainen atomi on täsmälleen hilarakenteen määrittelemässä tasapainokohdassa} 
	
	\item Oikeat kiteet eivät ole täydellisiä\appear{, vaan niissä on kidevirheitä:} \appear{\CDAlert[red]{puuttuvia atomeita}}\appear{, \CDAlert[red]{väärissä paikoissa olevia atomeita}}\appear{, \CDAlert[red]{epäsäännöllisyyksiä} atomien välisissä etäisyyksissä}\appear{, ja \CDAlert[red]{epäpuhtauksia}}
	
	\item \CDAlert[red]{Kun kidettä jännitetään} se muuttaa muotoaan \appear{(se \CDAlert[tunired]{venyy})} \appear{peruuttamattomasti kun rakenteelliset \CDAlert[red]{dislokaatiot} siirtävät atomeita}
	
\item Oikeat materiaalit \appear{(eivät ultrapuhtaat)} \appear{venyvät helpommin kuin mitä yksinkertainen teoria ennustaa}\appear{, mikä johtuu aineissa olevista kidevirheistä}\appear{, kuten särmäsiirtymistä/-dislokaatioista}

%Initial configuration of a crystal with an edge dislocation
\item Tarkastellaan pikaisesti miten tämä tapahtuu  \appear{(jättäen tarkemman tarkastelun kiinteän olomuodon fysiikan kurssillemme...)}
%\item The dislocation moves to the right as the atoms in the layer under it successively shift their bonds with those of the upper layer, one line at the time
	
\end{itemize}


\endofslide 
%\endofsection
\end{frame}
%%%%%%%%%%%%%%%


%%%%%%%%%%%%%%%
\begin{frame}
\frametitle{\texorpdfstring{\culine[\sectiontitleulinecolor]{\sectiontitle}}{\sectiontitle} }
\framesubtitle{Kidevirheet}
%Crystal defects
\appear{}

\begin{itemize}
	
%Initial configuration of a crystal with an edge dislocation
\item \CDAlert[fuchsia]{Dislokaatio} (siirros) on äkillinen muutos atomien järjestäytymisessä \appear{(Kuva)}%
\appear{\xdef\loc{south}%
\begin{tikzpicture}[remember picture,overlay]% anchor=south
    \node (A) [yshift=-0.05*\figYshift,xshift=-3.5*\figXshift, anchor=\loc] at (current page.\loc) {\includegraphics[page=1,height=5.5cm]{Figures_booklet/SOM_10_Bonding_figures/Tikz_SOM_Bonding_Mechanical_edge_dislocation_fixed}};%Tikz_SOM_Bonding_Mechanical_dislocation_movement
\end{tikzpicture}}%
\item Siirrokset \appear{(ja muut kidevirheet)} \appear{\Alert[fuchsia]{selittävät miksi oikeat materiaalit hajoavat ja muovautuvat} paljon helpommin} \appear{kuin mitä yksinkertainen (mutta järkevä) kiteisiin aineisiin perustuva malli ennustaa}

\end{itemize}

\endofslide 
%\endofsection
\end{frame}
%%%%%%%%%%%%%%%

%%%%%%%%%%%%%%%
\begin{frame}
\frametitle{\texorpdfstring{\culine[\sectiontitleulinecolor]{\sectiontitle}}{\sectiontitle} }
\framesubtitle{Kidevirheet}
%Crystal defects
\appear{}

\begin{itemize}
	
\item Dislokaatio liikkuu kuvissa oikealle kun sen lähellä olevien atomien sidokset hajoavat ja muodostuvat (toisten atomien kanssa) uudelleen\appear{, \Alert[fuchsia]{yksi atomirivi kerrallaan}} \appear{(Kuvat a–c)}
\appear{\xdef\loc{south}%
\begin{tikzpicture}[remember picture,overlay]% anchor=south
    \node (A) [yshift=-0.05*\figYshift,xshift=-0*\figXshift, anchor=\loc] at (current page.\loc) {\includegraphics[page=1,width=14cm]{Figures_booklet/SOM_10_Bonding_figures/Tikz_SOM_Bonding_Mechanical_dislocation_movement_fixed}};%Tikz_SOM_Bonding_Mechanical_dislocation_movement
\end{tikzpicture}}
\item Prosessi vaatii huomattavasti vähemmän energiaa tietyllä ajanhetkellä\appear{, kuin mitä tarvittaisiin useiden (tuhansien) sidosten yhtäaikaiseen uudelleenjärjestymiseen}

\item Tällaisia \CDAlert[fuchsia]{(plastisia) epämuodostumia} voidaan kontrolloida\appear{, synnyttämällä materiaaliin tarkoituksella kidevirheitä (Kuva d) }
	
\end{itemize}



\endofslide 
%\endofsection
\end{frame}
%%%%%%%%%%%%%%%

%%%%%%%%%%%%%%%
\begin{frame}
\frametitle{\texorpdfstring{\culine[\sectiontitleulinecolor]{\sectiontitle}}{\sectiontitle} }
\framesubtitle{Kiteen hilarakenne ja hilapisteet }
\appear{}

\begin{itemize}[itemsep=1.6mm]
	\item Kiinteät ja kiteiset aineet \Alert[red]{jaotellaan niiden kiteiden omaavien symmetrioiden perusteella} 7 kidejärjestelmään ja \Alert[red]{14 hilatyyppiin} 
	
	\item Kiteiden rakenteet\appear{, eli \CDAlert[fuchsia]{kidehilat}}\appear{, voidaan muodostaa toistamalla yksikkökoppia }
	
	\item Yksikkökoppi määritellään \CDAlert[blue]{kolmen translaatiovektorin $\mathbf{a}_i$ avulla}\appear{, missä $i=1$, 2, 3}\appear{, jolloin kaksi toisistaan erillään olevaa mutta keskenään ekvivalenttia \CDAlert[fuchsia]{hilapistettä} $\mathbf{r}^{\prime}$} \appear{ja $\mathbf{r}$}\appear{, liittyvät toisiinsa relaation kautta:}
\[ 
 \mathbf{r}^{\prime}  \equal \appear{\mathbf{r}}  \plus \appear{m_{1} \mathbf{a}_1} \plus \appear{m_{2} \mathbf{a}_2}  \plus \appear{m_{3} \mathbf{a}_3} \appear{\,, \qquad \text{missä}~m_{i}}\appear{=kokonaisluku} 
 \]

\item Yksikkökopit eivät ole yksikäsitteisiä\appear{, vaan kidehila voidaan muodostaa monen erilaisen yksikkökopin avulla}

\item \Alert[red]{Pienintä yksikkökoppia kutsutaan primitiiviseksi kopiksi/\CDAlert[red]{alkeiskopiksi}}\appear{, ja sen särmiä kutsutaan \CDAlert[tunired]{hilavakioiksi}}
%\item placed in each lattice point in same direction and posture
\end{itemize}

\endofslide 
%\endofsection
\end{frame}
%%%%%%%%%%%%%%%

%%%%%%%%%%%%%%%
\begin{frame}
\frametitle{\texorpdfstring{\culine[\sectiontitleulinecolor]{\sectiontitle}}{\sectiontitle} }
\framesubtitle{Kiderakenne, kuutiolliset kidejärjestelmät}
\appear{}

%\begin{itemize}
%	\item \CDAlert[red]{Body-centered cubic (bcc)}:
%\item[] \begin{itemize}
%	\item A repetition of cubes with atoms placed at the corners and at the center of the cube
%\item Solids from the periodic table's first column: Li, Na, K, Rb, and Cs
%\item Transition metals, such as: Fe, Cr, W
%\end{itemize}
%\item \CDAlert[blue]{Face-centered cubic (fcc)}:
%\item[] \begin{itemize}
%	\item A repetitive cubic structure, with atoms placed both at the corners of the cube and at the centers of the cubic faces
%	\item Noble gases (except helium) 
%	\item Transition metals such as: Cu, Ag, Au, Ni, Pb
%\end{itemize}
%\end{itemize}

%pintakeskinen kuutiollinen hila
\vspace{1mm}
\begin{minipage}{10.5cm}
\begin{itemize}%[itemsep=1.6mm] 
\item \CDAlert[red]{Tilakeskinen kuutiollinen hila (bcc)}:
\begin{itemize}
	\item On kuutiollinen kidehila\appear{, jossa kuution kärjissä} \appear{ja kuution keskellä on atomi}\appear{/\CDAlert[fuchsia]{hilapiste}}	
	\item Alkuainetaulukon 1 ryhmän kiinteät aineet: \appear{Li}\appear{, Na}\appear{, K}\appear{, Rb, ja Cs}
	\item Siirtymämetallit\appear{, kuten: Fe, Cr, W}
	\end{itemize}

\item \CDAlert[blue]{Pintakeskinen kuutiollinen hila (fcc)}:
\begin{itemize}
	\item Kuutiollinen kidehila\appear{, jossa kuution kärjissä} \appear{sekä tahkojen keskellä on atomit/\CDAlert[fuchsia]{hilapisteet}} 
	\item Jalokaasut \appear{(paitsi helium)}
	\item Siirtymämetallit \appear{kuten: Cu, Ag, Au, Ni, Pb}
	\end{itemize}
\end{itemize}
\end{minipage}
%% 
\xdef\loc{south east}
\begin{tikzpicture}[remember picture,overlay] % anchor=south
    \node (A) [yshift=-0.1*\figYshift,xshift=\figXshift, anchor=\loc] at (current page.\loc) {\includegraphics[page=1,height=9cm]{Figures_booklet/SOM_10_Bonding_figures/SOM_Bonding_Figure_22.png}}; %scale=\figscale. % 8cm max height
\end{tikzpicture}
\footnoteleft{\HarrisCR[1-]}

\endofslide 
%\endofsection
\end{frame}
%%%%%%%%%%%%%%%

%%%%%%%%%%%%%%
\begin{frame}
\frametitle{\texorpdfstring{\culine[\sectiontitleulinecolor]{\sectiontitle}}{\sectiontitle} }
\framesubtitle{Kiderakenteet, kuusikulmainen suljettu pakkaus}
\appear{}

\vspace{1mm}
\begin{minipage}{10.5cm}
\begin{itemize}%[itemsep=1.6mm] 
\item \CDAlert[red]{Kuusikulmainen suljettu pakkaus (hcp)}:
\begin{itemize}
	\item Kuusinkertainen kiertosymmetria
\item Alkuaineita joka puolelta alkuainetaulukkoa:
\item He, Be, Cd, Ce, Mg, Os, Zn, Zr\appear{, ja monet harvinaiset maametallit}

\end{itemize}

\item \CDAlert[red]{Yksinkertainen kuutiollinen hila (sc)}:
\begin{itemize}
	\item Toistuva kuutiollinen hila, jossa atomit/\CDAlert[fuchsia]{hilapisteet} ovat kuution kärjissä 
\end{itemize}

\vspace{3mm}

\item Huomaa että usein oletetaan: \appear{\CDAlert[fuchsia]{hilapiste = atomi}}

\item Hilapiste voi kuitenkin sisältää monta ato\-mia/mo\-le\-kyyliä\appear{, minkä johdosta kolmiulotteisia rakenteita on usein \CDAlert[blue]{vaikea visualisoida}}
\end{itemize}
\end{minipage}

%% 
\xdef\loc{south east}
\begin{tikzpicture}[remember picture,overlay] % anchor=south
    \node (A) [yshift=-0.1*\figYshift,xshift=\figXshift, anchor=\loc] at (current page.\loc) {\includegraphics[page=1,height=9cm]{Figures_booklet/SOM_10_Bonding_figures/SOM_Bonding_Figure_22.png}}; %scale=\figscale. % 8cm max height
\end{tikzpicture}
\footnoteleft{\HarrisCR[1-]}


\endofslide 
%\endofsection
\end{frame}
%%%%%%%%%%%%%%%

%%%%%%%%%%%%%%%
\begin{frame}
\frametitle{\texorpdfstring{\culine[\sectiontitleulinecolor]{\sectiontitle}}{\sectiontitle} }
\framesubtitle{3D pisteryhmät}
\appear{}

\begin{itemize}
	\item Kertauksena, kiteiset aineet muodostavat 7 kidejärjestelmää ja 14 hilaa, (\Alert[red]{linkki}{https://en.wikipedia.org/wiki/Space_group\#Table_of_space_groups_in_3_dimensions}):
	
%	\item Maybe only show the 2D groups here and mention that 3Ds will be covered in the future courses...?
\end{itemize}

%https://en.wikipedia.org/wiki/Space_group#Table_of_space_groups_in_3_dimensions
\xdef\loc{south}
\begin{tikzpicture}[remember picture,overlay] % anchor=south
    \node (A) [yshift=-0.1*\figYshift,xshift=0*\figXshift, anchor=\loc] at (current page.\loc) {\includegraphics[page=1,height=8.2cm]{Figures_booklet/SOM_10_Bonding_figures/Tikz_SOM_Bonding_Lattices_fixed}}; %scale=\figscale. % 8cm max height
\end{tikzpicture}

\endofslide 
%\endofsection
\end{frame}
%%%%%%%%%%%%%%%

%%%%%%%%%%%%%%%
\begin{frame}
\frametitle{\texorpdfstring{\culine[\sectiontitleulinecolor]{\sectiontitle}}{\sectiontitle} }
\framesubtitle{Kovalenttiset kiinteät aineet}
\appear{}

\begin{itemize}
	\item \CDAlert[red]{Kovalenttisesti sitoutuneiden aineiden} sidokset ovat samantyyppisiä kuin kovalenttisten molekyylien sidokset\appear{, atomit jakavat elektroneita}

\item Kovalenttisissa aineissa atomit muodostavat kovalenttisia sidoksia naapuriatomien kanssa 
 \implies{ yhtenäinen vahvojen sidosten verkosto}
 
 \item \Alert[red]{Tällaiset kiinteät aineet ovat kovia}\appear{, ja niillä on korkea sulamispiste}\appear{, ne ovat huonoja sähkönjohteita} \appear{(\CDAlert[blue]{dielektrisiä aineita} tai \CDAlert[blue]{puolijohteita})} \appear{koska kaikki \CDAlert[tuniblue]{valenssielektronit ovat vangittuina atomien välisiin sidoksiin} }

\end{itemize}

\endofslide 
%\endofsection
\end{frame}
%%%%%%%%%%%%%%%


%%%%%%%%%%%%%%%
\begin{frame}
\frametitle{\texorpdfstring{\culine[\sectiontitleulinecolor]{\sectiontitle}}{\sectiontitle} }
\framesubtitle{Kovalenttiset kiinteät aineet}
\appear{}

\begin{itemize}
	
 \item \CDAlert[violet]{Timantti} \appear{on malliesimerkki kovalenttisesta kiinteästä aineesta}\appear{, ja koostuu $fcc$-hilasta} \appear{jossa on kaksi hiiliatomia jokaisessa hilapisteessä} \implies{ Rakenne on erikoinen ja tunnetaan nimellä \CDAlert[violet]{timanttirakenne} }

\item Toisin sanoen\appear{ timantti koostuu kahdesta osittain päällekkäisestä fcc-hilasta} 

\item Tiivis rakenne johtaa myös suurimpaan atomitason koheesioenergiaan mitä kiinteille ainelle löytyy\appear{, $7,37 \mathrm{~eV}$/atomi}

\end{itemize}

\vspace{2.5mm}
\begin{minipage}{8cm}
\begin{itemize}%[itemsep=1.6mm] 
	\item Huomaa että hiilellä on muitakin tunnettuja kiderakenteita\appear{, kuten \CDAlert[tunired]{grafiitti}}\appear{, kiinteät \CDAlert[fuchsia]{fullereneenit}} \appear{ja \CDAlert[red]{grafeeni}}
\end{itemize}
\end{minipage}

\xdef\loc{south east}
\begin{tikzpicture}[remember picture,overlay] % anchor=south
    \node (A) [yshift=-0.1*\figYshift,xshift=\figXshift, anchor=\loc] at (current page.\loc) {\includegraphics[page=1,height=5cm]{Figures_booklet/SOM_10_Bonding_figures/Tikz_SOM_Bonding_diamond_fixed}}; %scale=\figscale. % 8cm max height
\end{tikzpicture}

\endofslide 
%\endofsection
\end{frame}
%%%%%%%%%%%%%%%


%%%%%%%%%%%%%%%
\begin{frame}
\frametitle{\texorpdfstring{\culine[\sectiontitleulinecolor]{\sectiontitle}}{\sectiontitle} }
\framesubtitle{Joitakin hiilen monista allotroopeista}
\appear{}

\xdef\loc{south west}
\begin{tikzpicture}[remember picture,overlay] % anchor=south
    \node (A) [yshift=-0.1*\figYshift,xshift=-1*\figXshift, anchor=\loc] at (current page.\loc) {\includegraphics[page=1,width=5cm]{Figures_booklet/SOM_10_Bonding_figures/Tikz_SOM_Bonding_figures_graphene_graphite_fixed}}; %scale=\figscale. % 8cm max height
\end{tikzpicture}

\xdef\loc{west}
\begin{tikzpicture}[remember picture,overlay] % anchor=south
    \node (A) [yshift=-1.5*\figYshift,xshift=-1*\figXshift, anchor=\loc] at (current page.\loc) {\includegraphics[page=2,width=5.5cm]{Figures_booklet/SOM_10_Bonding_figures/Tikz_SOM_Bonding_figures_graphene_graphite_fixed}}; %scale=\figscale. % 8cm max height
\end{tikzpicture}
%Tikz_SOM_Bonding_figures_graphene_graphite

\xdef\loc{south}
\begin{tikzpicture}[remember picture,overlay] % anchor=south
    \node (A) [yshift=-0.1*\figYshift,xshift=0*\figXshift, anchor=\loc] at (current page.\loc) {\includegraphics[page=1,height=8.5cm]{Figures_booklet/SOM_10_Bonding_figures/SOM_Bonding_Figure_58.png}}; %scale=\figscale. % 8cm max height
\end{tikzpicture}

\xdef\loc{south east}
\begin{tikzpicture}[remember picture,overlay] % anchor=south
    \node (A) [yshift=-0.35*\figYshift,xshift=\figXshift, anchor=\loc] at (current page.\loc) {\includegraphics[page=1,height=6cm]{Figures_booklet/SOM_10_Bonding_figures/SOM_Bonding_Figure_57.png}}; %scale=\figscale. % 8cm max height
\end{tikzpicture}
\footnote{\HarrisCRshort[1-]}

\endofslide 
%\endofsection
\end{frame}
%%%%%%%%%%%%%%%

%%%%%%%%%%%%%%%
\begin{frame}
\frametitle{\texorpdfstring{\culine[\sectiontitleulinecolor]{\sectiontitle}}{\sectiontitle} }
\framesubtitle{Ioniset kiinteät aineet}
\appear{}

\begin{itemize}
	\item Ruokasuolan $\mathrm{NaCl}$\appear{ ionit Na$^+$ ja Cl$^-$ ovat keskeissymmetrisiä}\appear{, mutta Cl$^-$-ioni on lähes kaksi kertaa niin suuri kuin Na$^+$}

\item Potentiaalienergian minimi saavutetaan kun yhdellä ionilla (Na$^+$ tai Cl$^-$) on kuusi toisentyyppistä ionia \CDAlert[blue]{lähimpinä naapureinaan}  
\implies{ $\mathrm{NaCl}$ muodostuu pintakeskeisestä kuutiollisesta hilasta (fcc)\appear{, jonka \CDAlert[red]{hilapisteet} \CDAlert[red]{koostuvat yhdestä Na ja yhdestä Cl atomista}}} 
\end{itemize}

\vspace{2.5mm}
\begin{minipage}{8cm}
\begin{itemize}%[itemsep=1.6mm] 
	\item Voisi myös ajatella että NaCl koostuu kahdesta sisäkkäisestä fcc-hilasta \appear{(toinen Na:lle}\appear{, toinen Cl:lle)}
\item On hyvä tietää että ruokasuola (NaCl) on todella tärkeä kiin\-teä aine:\\ \appear{\CDAlert[tunired]{jokainen elävä organismi}} \appear{\Alert[tunired]{Maassa tarvitsee suolaa}}

\end{itemize}
\end{minipage}

\xdef\loc{south east}
\begin{tikzpicture}[remember picture,overlay] % anchor=south
    \node (A) [yshift=-0.25*\figYshift,xshift=\figXshift, anchor=\loc] at (current page.\loc) {\includegraphics[page=1,height=5cm]{Figures_booklet/SOM_10_Bonding_figures/Tikz_SOM_Bonding_NaCl_fixed}}; %scale=\figscale. % 8cm max height
\end{tikzpicture}

\endofslide 
%\endofsection
\end{frame}
%%%%%%%%%%%%%%%

%%%%%%%%%%%%%%%
\begin{frame}
\frametitle{\texorpdfstring{\culine[\sectiontitleulinecolor]{\sectiontitle}}{\sectiontitle} }
\framesubtitle{Ioniset kiinteät aineet ja Madelungin vakio}
\appear{}

\begin{itemize}
	\item Huomaa että $\mathrm{Na}$ ja $\mathrm{Cl}$ eivät muodosta molekyylejä $\mathrm{NaCl}$ kiinteässä rakenteessa 
	
	\item Ioniparin sähköinen energia voidaan kuitenkin laskea:
\[ 
 U_{\text {sitova}}(r)  \equal \appear{-\alpha} \frac{1}{4 \pi \varepsilon_{0}} \frac{e^{2}}{r}  \equal \appear{-\alpha} \frac{k e^{2}}{r} \appear{\,,\qquad \text{''} \Frac{1}{4 \pi \varepsilon_{0}} \Equiv k \text{''} } 
 \]
\appear{missä  $r$ on vierekkäisten ionien välinen etäisyys } \appear{$(0,282 \mathrm{~nm}~\mathrm{NaCl}$:lle)}

\item Vakiota $\alpha$ kutsutaan \CDAlert[red]{Madelungin vakioksi}\appear{, joka ottaa huomioon naapurien aiheuttamat sähköstaattiset hylkimiset}\appear{/puoleensävetämiset} 

\item Jos vain 6 lähintä vastakkaisesti varautunutta naapuria vaikuttaisi Madelungin vakion $\alpha$ suuruuteen\appear{, sen arvoksi tulisi 6} 

\end{itemize}

\endofslide 
%\endofsection
\end{frame}
%%%%%%%%%%%%%%%


%%%%%%%%%%%%%%%
\begin{frame}
\frametitle{\texorpdfstring{\culine[\sectiontitleulinecolor]{\sectiontitle}}{\sectiontitle} }
\framesubtitle{Madelungin vakio}
\appear{}

\begin{itemize}
	
\item Mutta kiderakenteessa on myös 12 samanvarauksista ionia etäisyyden $2^{1 / 2} r$ päässä\appear{, ja 8 vastakkaismerkkistä ionia etäisyyksillä $3^{1 / 2} r$ jne.} \appear{Eli yhteisvaikutukseksi saadaan:}
\[ 
 \alpha  \equal \appear{6-} \frac{12}{\sqrt{2}} \plus \frac{8}{\sqrt{3}}  \minus \frac{6}{2} \plus \frac{20}{\sqrt{5}}  \minus \appear{\ldots}
 \]
\appear{mikä johtaa summaan joka ei konvergoidu!} 

\item Ryhmittelemällä termejä fiksusti\appear{, summa saadaan konvergoitumaan}\appear{, hitaasti, ja tulokseksi saadaan $\alpha=1,7476$}
\end{itemize}

\endofslide 
%\endofsection
\end{frame}
%%%%%%%%%%%%%%%

%%%%%%%%%%%%%%%
\begin{frame}
\frametitle{\texorpdfstring{\culine[\sectiontitleulinecolor]{\sectiontitle}}{\sectiontitle} }
\framesubtitle{Madelungin vakio}
\appear{}

\begin{itemize}
	\item Taulukossa on annettu joidenkin yleisten kiinteiden aineiden ominaisuuksia
\end{itemize}
\appear{\xdef\loc{south east}%
\begin{tikzpicture}[remember picture,overlay] % anchor=south
    \node (A) [yshift=-0.1*\figYshift,xshift=\figXshift, anchor=\loc] at (current page.\loc) {\includegraphics[page=1,height=8.2cm]{Figures_booklet/SOM_Tipler_Solid_state_physics_figures/SOM_Tipler_SSP_figure_5.png}};%scale=\figscale. % 8cm max height
\end{tikzpicture}\footnoteleft{\TiplerCRshort}
}

\vspace{2.5mm}
\begin{minipage}{5.5cm}
\begin{itemize}%[itemsep=1.6mm] 
\item Huomaa että Madelungin vakio soveltuu vain ionisille kiinteille aineille 

\item Ja että vakion $\alpha$ arvo riippuu aineen kiderakenteesta
\end{itemize}
\end{minipage}


\endofslide 
%\endofsection
\end{frame}
%%%%%%%%%%%%%%%

%%%%%%%%%%%%%%%
\begin{frame}
\frametitle{\texorpdfstring{\culine[\sectiontitleulinecolor]{\sectiontitle}}{\sectiontitle} }
\framesubtitle{Ioniset kiinteät aineet}
\appear{}

\begin{itemize}[itemsep=1.8mm]
	\item Kun ionit ovat lähellä toisiaan\appear{, elektronien aaltofunktiot alkavat mennä päällekkäin} \appear{ja \CDAlert[red]{Paulin kieltosääntö}} \appear{synnyttää ionien välille hylkivän voiman}
	
	\item Tässä tapauksessa potentiaalienergian oletetaan olevan muotoa
\[ 
 U(r)  \equal \appear{-\alpha} \frac{k e^{2}}{r}  \plus \frac{A}{r^{n}} 
 \]

\item Hakemalla energian derivaatan nollakohta ($r$:n suhteen)\appear{, tasapainoetäisyydellä $r=r_{0}$}\appear{ saadaan vakioksi $A$ ratkaistua}
\[ 
 A  \equal  \frac{\alpha k e^{2} r_{0}^{n-1}}{n} 
 \]

\item Jolloin potentiaalienergia saadaan muotoon
\[ 
 U(r)  \equal \appear{-\alpha} \frac{k e^{2}}{r_{0}} \LB \frac{r_{0}}{r} \appear{-\Frac{1}{n}} \appear{\bigg(} \frac{r_{0}}{r} \appear{\bigg)^{n}} \;\RB
 \]
% \appear{\textrm{[[}  \appear{\textrm{]]}}
\end{itemize}

\endofslide 
%\endofsection
\end{frame}
%%%%%%%%%%%%%%%

%%%%%%%%%%%%%%%
\begin{frame}
\frametitle{\texorpdfstring{\culine[\sectiontitleulinecolor]{\sectiontitle}}{\sectiontitle} }
\framesubtitle{Ioniset kiinteät aineet}
\appear{}

\begin{itemize}[itemsep=1.8mm]

\item Joka saadaan tasapainopisteessä ($r=r_{0}$) muotoon\vspace{-1.5mm}

\[ 
 U \textrm{((}r_{0} \textrm{))}  \equal \appear{-\alpha} \frac{k e^{2}}{r_{0}} \appear{\bigg[} \appear{1-} \frac{1}{n} \appear{\bigg]}
 \]

	\item Tasapainoetäisyys voidaan määrittää kokeellisesti esim. \CDAlert[red]{röntgendiffraktiolla} \appear{(tai \CDAlert[blue]{laskemalla se kiteen tiheydestä}) }%
	\vspace{-0.5mm}
	\implies{ Ionisen kiinteän aineen dissosiaatio- tai hilaenergia kertoo meille $n$ }
	\item  $\mathrm{NaCl}$:n dissosiaatioenergia on $770 \mathrm{\;kJ} / \mathrm{mol}\appear{=7,98$\;eV/\,ionipari}\appear{, johtaen $\appear{n=9,35} \approx \appear{9}$}
\end{itemize}

\vspace{1.6mm}
\begin{minipage}{8.8cm}
\begin{itemize}[itemsep=1.6mm]
	 
\item[] \begin{itemize}
	\item Tässä arvo $7,98$\;eV on energia joka tarvitaan poistamaan Na$^+$ ja Cl$^-$ ionipari kiteestä 
	\item $\mathrm{Cl}$:n elektroniaffiniteetti on $3,62$\;eV
	\item Na:n muodostaminen Na$^+$-ionista \appear{vapauttaa $5,15 \mathrm{\;eV}$ energiaa} 
\end{itemize}

\end{itemize}
\end{minipage}

\xdef\loc{south east}
\begin{tikzpicture}[remember picture,overlay] % anchor=south
    \node (A) [yshift=-0.35*\figYshift,xshift=\figXshift, anchor=\loc] at (current page.\loc) {\includegraphics[page=1,height=3.85cm]{Figures_booklet/SOM_10_Bonding_figures/SOM_Bonding_Figure_16.png}}; %scale=\figscale. % 8cm max height
\end{tikzpicture}
\footnote{\HarrisCR[1-]}

\endofslide 
%\endofsection
\end{frame}
%%%%%%%%%%%%%%%

%%%%%%%%%%%%%%%
\begin{frame}
\frametitle{\texorpdfstring{\culine[\sectiontitleulinecolor]{\sectiontitle}}{\sectiontitle} }
\framesubtitle{Ioniset kiinteät aineet}
\appear{}

\begin{itemize}

\item Yhden neutraalin $\mathrm{Na}$ \appear{ja $\mathrm{Cl}$ atomiparin} \appear{poistamiseksi kiteestä tarvitaan energiaa siis yhteensä} $\appear{7,98 \mathrm{\;eV}} \plus \appear{3,62 \mathrm{\;eV}}  \minus \appear{5,15 \;\mathrm{eV}} \equal \appear{6,45 \mathrm{\;eV}}$

\item Lopulta niin sanottu \CDAlert[red]{kiteen koheesioenergia}\appear{, mikä kerrotaan atomille tai atomiparille}\appear{, on $\Frac{6,45 \mathrm{\;eV}}{2}=3,225$~eV/atomi}\appear{, mikä on lähellä kirjallisuusarvoa $3,19 \mathrm{\;eV}/$atomi} 

\item \CDAlert[tunired]{Suuri koheesioenergia implikoi korkeaa sulamispistettä}

	\item $\mathrm{NaCl}$:n lisäksi\appear{, suurin osa ionisista kiteistä muodostaa fcc-hiloja:} \appear{LiF, KF, KCl, KI, AgCl, ...}

\item Toinen yleinen ionikiteiden rakenne on yksinkertainen kuutiollinen hila (sc)\appear{, joita muodostavat CsCl}\appear{, TlI, TlBr, LiHg, NH$_{4}$,  Cl, ...}

\item Madelungin vakio sc-hilalle on $\alpha=1,7627$ 
\end{itemize}


\endofslide 
%\endofsection
\end{frame}
%%%%%%%%%%%%%%%

%%%%%%%%%%%%%%%
\begin{frame}
\frametitle{\texorpdfstring{\culine[\sectiontitleulinecolor]{\sectiontitle}}{\sectiontitle} }
\framesubtitle{Ioniset kiinteät aineet ja Madelungin vakion esimerkki}
\appear{}

\begin{itemize}
	\item \CDAlert[red]{Esimerkki}: Lasketaan Madelungin vakio kuvan kaksiulotteiselle kiteelle%
	\appear{\xdef\loc{south east}%
\begin{tikzpicture}[remember picture,overlay]% anchor=south
    \node (A) [yshift=-0.15*\figYshift,xshift=\figXshift, anchor=\loc] at (current page.\loc) {\includegraphics[page=1,height=6cm]{Figures_booklet/SOM_Tipler_Solid_state_physics_figures/SOM_Tipler_SSP_figure_9.png}};%scale=\figscale. % 8cm max height
\end{tikzpicture}\footnoteleft{\TiplerCR}
}
 [Vinkki: lasketaan vain 4 ensimmäistä termiä]
	
\end{itemize}

\vspace{1.6mm}
\begin{minipage}{6.5cm}
\begin{itemize}[itemsep=1.6mm]
	 
\item[] \begin{itemize}
	\item[] $$
\begin{aligned}
&\appear{U_{\text {vetävä}}} \appear{=}\appear{-k e^{2}} \LP \frac{4}{r} \minus \frac{4}{2^{1/2} r} \minus \frac{4}{2 r} \plus \frac{8}{5^{1/2} r} \,\RP \\
&\appear{U_{\text {vetävä}}}  \equal  \minus \frac{k e^{2}}{r} \LP \appear{4} \minus \frac{4}{2^{1/2}} \minus \appear{2} \plus \frac{8}{5^{1/2}} \,\RP
\end{aligned}
$$
\item[] Vastaus: \appear{$\alpha \Approx 2,749$} 
\end{itemize}

\end{itemize}
\end{minipage}

\endofslide 
%\endofsection
\end{frame}
%%%%%%%%%%%%%%%

%%%%%%%%%%%%%%%
\begin{frame}
\frametitle{\texorpdfstring{\culine[\sectiontitleulinecolor]{\sectiontitle}}{\sectiontitle} }
\framesubtitle{Metallit}
\appear{}

\begin{itemize}[itemsep=1.55mm]
	\item Paitsi jalokaasut, kaikilla alkuaineilla on valenssielektroneja 
	
	\item Suurimmassa osassa alkuaineista \appear{atomin kaikki valenssielektronit eivät voi muodostaa kovalenttisia sidoksia naapuriatomien kanssa} 	
%	But in most cases no arrangement of covalent bonds can join all of a given atom's valence electrons to those of the atoms surrounding it 
	
	\item Alkuaineet tai yhdisteet  \CDAlert[tunired]{jolla on ylimääräisiä valenssielektroneja } \appear{muodostavat \CDAlert[red]{metalleja}}

\item Käytännössä muodostuu tilojen jatkumo\appear{, joiden energiat ovat yksittäisten atomien tiloja matalammalla }
\end{itemize}

\vspace{1.6mm}
\begin{minipage}{7.5cm}
\begin{itemize}[itemsep=1.5mm]
	 \item Koska valenssielektronit \appear{jakaantuvat kaikkien atomien kesken}\appear{, ne voivat liik\-kua lähestulkoon vapaasti hilarakenteessa}\appear{, muodostaen \CDAlert[red]{elektronikaasun}} 
\implies{ Niistä muodostuu \CDAlert[tunired]{johtavuuselektroneita} }

%\item \Alert[red]{Metallit ovat usein erinomaisia}\\ \CDAlert[red]{sähkönkohteita}
\item \Alert[red]{Metallit johtavat hyvin} \CDAlert[red]{sähköä ja}\\ \CDAlert[red]{lämpöä}

\end{itemize}
\end{minipage}


\xdef\loc{south east}
\begin{tikzpicture}[remember picture,overlay] % anchor=south
    \node (A) [yshift=-0.15*\figYshift,xshift=\figXshift, anchor=\loc] at (current page.\loc) {\includegraphics[page=1,height=4.4cm]{Figures_booklet/SOM_10_Bonding_figures/Tikz_SOM_Bonding_e_conduction_schematic_fixed}}; %scale=\figscale. % 8cm max height
\end{tikzpicture}

\endofslide 
%\endofsection
\end{frame}
%%%%%%%%%%%%%%%

%%%%%%%%%%%%%%%
\begin{frame}
\frametitle{\texorpdfstring{\culine[\sectiontitleulinecolor]{\sectiontitle}}{\sectiontitle} }
\framesubtitle{Metallit}
\appear{}

\begin{itemize}
	\item Tarkastellaan esimerkiksi litiumia:

\item[] \begin{itemize}
	\item Litiumin elektronirakenne on $\appear{1 s^{2}} \appear{2 s}$ 

\item Radiaalinen $2 s$-elektronin aaltofunktio kokee vetyatomimaisen ytimen\appear{, mikä koostuu atomin ytimestä ja täydestä 1s elektronikuoresta.}\appear{ $2 s$-elektronin aaltofunktio on muotoa}
\[ 
 \psi_{2,0} \propto \appear{R_{n=2, \ell=0}} \qquad \appear{\&} \qquad \appear{\psi_{2,0}} \equal \appear{C_{2,0}} \appear{\textrm{((}2} \minus \frac{r}{a_{0}} \appear{\textrm{))} }\appear{e^{-r / 2 a_{0}} }
 \]
\appear{missä $R_{n=2, \ell=0}$} \appear{on $2 s$-elektronin aaltofunktion radiaalinen riippuvuus} \appear{(koko aaltofunktio on separoituva)}\appear{, $C_{2,0}$ on normeerausvakio} \appear{$ \textrm{((}C_{2,0}= \textrm{((}2 a_{0} \textrm{))}^{-3 / 2} \textrm{))}$}\appear{, ja $a_{0}=$~Bohrin säde}

\end{itemize}

\end{itemize}

\endofslide 
%\endofsection
\end{frame}
%%%%%%%%%%%%%%%

%%%%%%%%%%%%%%%
\begin{frame}
\frametitle{\texorpdfstring{\culine[\sectiontitleulinecolor]{\sectiontitle}}{\sectiontitle} }
\framesubtitle{Metallit}
\appear{}

\begin{itemize}
	\item Tarkastellaan esimerkiksi litiumia:

\item[] \begin{itemize}
	\item Kuvassa on hahmoteltu todennäköisyystiheyksiä:
	\item[] \begin{itemize}
			\item[(a)] yksittäiselle litium-atomille
			\item[(b)] metalliselle litiumille
		\end{itemize}
\end{itemize}
\end{itemize}

\vspace{1.6mm}
\begin{minipage}{4.5cm}
\begin{itemize}[itemsep=1.6mm]
\item[] \begin{itemize}
	 \item Todennäköisyys\-jakau\-ma lähellä Li-atomeita ($\pm m\cdot0,3 \mathrm{\;nm}$, $m=\mathbb{Z}$) kaventuu\appear{, mutta elektroni on silti hyvin \Alert[blue]{delokalisoitunut laajalle alueelle}} %and spatial locations of the electrons are more peaked around 

\end{itemize}

	
\end{itemize}
\end{minipage}

\xdef\loc{south east}
\begin{tikzpicture}[remember picture,overlay] % anchor=south
    \node (A) [yshift=-0.15*\figYshift,xshift=\figXshift, anchor=\loc] at (current page.\loc) {\includegraphics[page=1,height=6cm]{Figures_booklet/SOM_Tipler_Solid_state_physics_figures/SOM_Tipler_SSP_figure_11.png}}; %scale=\figscale. % 8cm max height
\end{tikzpicture}
\footnoteleft{\TiplerCR[1-]}

\endofslide 
%\endofsection
\end{frame}
%%%%%%%%%%%%%%%

%%%%%%%%%%%%%%%
\begin{frame}
\frametitle{\texorpdfstring{\culine[\sectiontitleulinecolor]{\sectiontitle}}{\sectiontitle} }
\framesubtitle{Molekylääriset kiinteät aineet}
\appear{}

\begin{itemize}
	\item Molekylaarisissa kiinteissä aineissa \appear{hilapisteissä} \appear{on atomien sijasta} \appear{molekyylejä} \appear{(vertaa NaCl)}
	
	\item Koheesiovoimat jotka sitovat molekyylit toisiinsa ovat \CDAlert[red]{Van der Waalsin voimia}\appear{: \CDAlert[tunired]{dipoli–dipoli vuorovaikutuksia}}\appear{, \CDAlert[red]{Londonin voimia} jne.}
	
	\item Vaikka kovalenttisen molekyylin elektronit muodostavat stabiilin molekyylin\appear{, eikä ole energeettisesti järkevää jakaa elektroneita useamman molekyylin kesken}\appear{, molekyylien välillä on silti puoleensavetäviä voimia}\appear{, kun molekyylit ovat tarpeeksi lähellä toisiaan}

\item Symmetriset kovalenttiset molekyylit \appear{(esim. $\mathrm{N}_{2}$ tai $\mathrm{H}_{2}$)} \appear{joilla ei ole pysyvää sähködipolimomenttia}\appear{, voivat silti polarisoitua} \appear{(sähködipolimomentti indusoituu molekyyliin)}  
\implies{ \CDAlert[red]{indusoitu dipoli–dipoli vuorovaikutus} tunnetaan nimellä \CDAlert[red]{Londonin voima} }

\end{itemize}

\endofslide 
%\endofsection
\end{frame}
%%%%%%%%%%%%%%%

%%%%%%%%%%%%%%%
\begin{frame}
\frametitle{\texorpdfstring{\culine[\sectiontitleulinecolor]{\sectiontitle}}{\sectiontitle} }
\framesubtitle{Molekylääriset kiinteät aineet}
\appear{}

\begin{itemize}
	\item Jos molekyyleillä on pysyvä sähködipolimomentti \appear{(esim. $\mathrm{NH}_{3}$ ja $\mathrm{H}_{2} \mathrm{O}$)}\appear{, niin vielä voimakkaampi}\appear{ (eli tärkeämpi)}\appear{, \CDAlert[red]{dipoli–dipoli vuorovaikutus} muodostuu } 
	\implies{ Tämä on erityisen voimakasta  \CDAlert[blue]{vedessä}, missä vuorovaikutusta\\
	 kutsutaan \CDAlert[blue]{vetysidokseksi} }

\item \CDAlert[red]{Molekylääriset kiinteät aineet}:
\item[] \begin{itemize}
	\item Ovat suhteellisen pehmeitä 
	\item Omaavat matalan sulamispisteen \appear{$ \textrm{((}\!<200^{\circ}C \textrm{))}$}
\item Ovat huonoja lämmönjohteita \appear{ja sähkönjohteita}
\item \CDAlert[blue]{Vesijää $H_2 O$} \appear{ja kuivajää (eli hiilihappojää) $CO_2$} \appear{ovat malliesimerkkejä tällaisista kiinteistä aineista}
\end{itemize}

\end{itemize}

\endofslide 
%\endofsection
\end{frame}
%%%%%%%%%%%%%%%

%%%%%%%%%%%%%%%
\begin{frame}
\frametitle{\texorpdfstring{\culine[\sectiontitleulinecolor]{\sectiontitle}}{\sectiontitle} }
\framesubtitle{Yhteenveto}
\appear{}

\footnote{Encyclopaedia Britannica: metallic bond (chemistry)}
% https://www.britannica.com/science/metallic-bond
\xdef\loc{center}
\begin{tikzpicture}[remember picture,overlay] % anchor=south
    \node (A) [yshift=0.35*\figYshift,xshift=0*\figXshift, anchor=\loc] at (current page.\loc) {\includegraphics[page=1,width=14cm]{Figures_booklet/SOM_bonding_types_of_bonds.png}}; %scale=\figscale. % 8cm max height
\end{tikzpicture}

%\endofslide 
%\endofsection
\end{frame}
%%%%%%%%%%%%%%%